\documentclass[instructions.tex]{subfiles}
\begin{document}
 
\chapter{De l’amour du prochain.}
 
Cette vertu est nécessaire : quels en sont les motifs ? Cette vertu est rare : quels en sont les obstacles ?
 
\primo Vertu si nécessaire, que si entre tous les hommes qui sont sur la terre il y en avoit un seul que nous n’aimassions pas, ce seroit assez pour être rejetés de Dieu. Quand je parlerois le langage des anges, dit saint Paul, si je n’ai pas la charité, je ne suis rien. Voilà le grand commandement du Sauveur, si vous l’accomplissez, dit saint Jean, tout est accompli. Vertu si excellente, qu’elle est au-dessus de tous les sacrifices\footnote{Magis est omnibus holocaustis et sacrificiis. Marc. 12.}. Vertu aussi indispensable que l’amour de Dieu ; puisque le précepte d’aimer le prochain est semblable à celui d’aimer Dieu, et qu’on ne peut accomplir l’un sans l’autre. Celui qui hait son frère, et qui dit qu’il aime Dieu, est un menteur, dit saint Jean ; car s’il n’aime pas son prochain, qu’il voit, comment aimera-t-il Dieu, qu’il ne voit pas\footnote{Non sum sicut cæteri. Luc. 18.}.
 
Vous devez donc aimer le prochain, mais par quels motifs ? Vous devez l’aimer, parce que Dieu commande de l’aimer, et parce que Dieu l’aime. Vous devez l’aimer, parce qu’il est votre frère, l’image et l’enfant de Dieu. Vous devez l’aimer, parce qu’il est aimé de Jésus-Christ, qu’il est racheté de son sang, qu’il est destiné au ciel aussi bien que vous. Vous seriez bien déraisonnable de ne pas aimer celui qu’un Dieu a plus chéri que sa propre vie. Vous devez enfin aimer le prochain, quel qu’il soit, parce que les sentiments que vous en aurez rejaillissent sur Jesus-Christ. Si vous haïssez, ou si vous aimez votre frère, c’est Jesus-Christ que vous haïssez, ou que vous aimez.
 
J’aime tout le monde, direz-vous, je n’ai point d’ennemis. Ce n’est pas assez de le dire. Aimez votre prochain, non pas en paroles, mais aimez-le comme vous-même, dit Jésus-Christ. Vous devez donc l’aimer plus que vos biens, et les perdre plutôt que de perdre la charité\footnote{Perde pecuniam propter fratrem et amicum. Eccl. 29.}. Vous devez faire aux autres ce que vous voudriez raisonnablement qu’on fît à vous-même; penser d’eux comme vous voudriez que l’on pensât de vous; ne jamais leur faire et ne jamais leur dire ce que vous ne voudriez pas qu’on vous fît ou qu’on dit de vous.
 
Un solitaire, ayant mené en apparence une vie fort imparfaite, se trouva néanmoins si consolé à la mort, que son supérieur lui ayant demandé d’où lui venoit ce grand contentement, il lui répondit : Mon père, j’ai tâché toute ma vie de pratiquer le grand précepte du Seigneur, d’aimer tout le monde, de supporter mes frères, de leur rendre service, de ne faire mal à personne, de ne parler, de ne juger mal de qui que ce soit, et de penser du bien de tous. Voilà ce qui fait ma consolation et qui, malgré mes imperfections, me donne tant de confiance pour aller paroître devant Dieu. — Ah ! mon frère, lui dit le supérieur, mourez en paix; heureux d’avoir vécu dans de si saintes dispositions !
 
\secundo Oh ! que cette vertu est rare et inconnue ! et que d’illusions sur ce point ! Tel croit pratiquer de grandes vertus qui ignore celle qui donne le prix à toutes les autres. Combien de personnes, sous les dehors d’une vie régulière, cachent une humeur dure et farouche, et qui ne savent ce que c’est que bienveillance, affabilité, ménagement, condescendance ! Tel passe pour un bon chrétien, qui n’a pas même les principes du christianisme; qui se choque de tout, qui s’irrite d’une parole, qui s’ombrage de rien; qui n’épargne, qui n’excuse, qui ne supporte personne. Y a-t-il manières plus sèches, ressentimens plus vifs, railleries plus mordantes, médisances plus artificieuses, que dans certaines gens qui se croient offensés, et qui se piquent cependant d’être irréprochables ?
 
Eh ! hypocrites, ignorez-vous qu’un acte de douceur et de charité vaut mieux que tous les sacrifices ? Semblables aux pharisiens, vous vous savez bon gré de n’être pas comme les autres. Mais ne savez-vous pas que sans la charité vous n’avez qu’un fantôme de religion, et que c’est par elle qu’on connoît les vrais serviteurs de Dieu ?
 
Les obstacles à la charité sont ordinairement ces trois vices : l’avarice, l’orgueil et l’envie. Vous n’aimez pas un homme, parce qu’il n’est pas dans vos intérêts ou parce qu’il vous a fait tort dans vos biens : voilà un effet de l’avarice. Vous ne pouvez souffrir une personne, parce qu’elle vous contredit et n’est pas de votre sentiment, ou parce qu’elle a dit une parole contre vous, ou parce qu’elle réussit et qu’elle prospère, ou parce qu’on l’approuve et qu’on l’estime, tandis qu’on ne pense pas à vous; c’est orgueil et envie : trois vices que Dieu a en horreur.
 
Vous regardez un scélérat comme un abominable; mais vous êtes encore plus abominable aux yeux de Dieu, si vous êtes sans charité. Un grand pécheur qui a des sentimens de charité pour tous, est bien près du royaume de Dieu, et sera plus tôt converti qu’un chrétien qui se croit vertueux, et qui vit sans charité.
 
Si votre prochain a de grands défauts, croyez que vous en avez encore de plus grands que vous ne connoissez pas; n’aimez pas ses défauts et ses vices, mais aimez sa personne. Les défauts d’autrui ne sont pas un titre pour le haïr, autrement vous haïriez tous les hommes, puisque aucun n’est exempt de défauts. Plus les défauts du prochain sont grands, plus il est digne de votre charité; supportez-le, ayez-en compassion; n’en parlez pas, à moins qu’il ne soit nécessaire ou utile, et priez pour lui.
 
\section{Résolutions}
 
\primo Afin d’aimer saintement mon prochain, je considérerai dans sa personne l’œuvre de Dieu, le prix du sang de Jésus-Christ; mon compagnon dans le voyage périlleux de cette vie, destiné, comme moi, au bonheur du ciel. 

\secundo Je me garderai bien de lui donner le moindre scandal

\tertio Si je remarque en lui des défauts, ou même des vices, 'en aurai compassion, je prierai pour lui, et je penserai qu'il vaut peut-être miueux que mois; que ce que j'apercois en lui de défectueux y est peut-être racheté par quelques vertu que je ne vois pa; que; s'il avait recu les mêmes grâces que moi, il serait peut-être moins imparfait que je ne le suis, et qu'enfin il sera peut-être beaucoup plus élevé dans le ciel. 

\prayer{O mon Dieu ! pardonnez-moi les fautes innombrables que j'ai faites jusqu'ici contre la charité; et daignez affermir tellement cette vertu dans mon coeur, que je ne la perde jamais plus} 

\end{document}
